% Options for packages loaded elsewhere
\PassOptionsToPackage{unicode}{hyperref}
\PassOptionsToPackage{hyphens}{url}
\PassOptionsToPackage{dvipsnames,svgnames,x11names}{xcolor}
\documentclass[
  10pt,
  fontset=fandol]{ctexart}
\usepackage{xcolor}
\usepackage[margin=16mm]{geometry}
\usepackage{amsmath,amssymb}
\setcounter{secnumdepth}{-\maxdimen} % remove section numbering
\usepackage{iftex}
\ifPDFTeX
  \usepackage[T1]{fontenc}
  \usepackage[utf8]{inputenc}
  \usepackage{textcomp} % provide euro and other symbols
\else % if luatex or xetex
  \usepackage{unicode-math} % this also loads fontspec
  \defaultfontfeatures{Scale=MatchLowercase}
  \defaultfontfeatures[\rmfamily]{Ligatures=TeX,Scale=1}
\fi
\usepackage{lmodern}
\ifPDFTeX\else
  % xetex/luatex font selection
\fi
% Use upquote if available, for straight quotes in verbatim environments
\IfFileExists{upquote.sty}{\usepackage{upquote}}{}
\IfFileExists{microtype.sty}{% use microtype if available
  \usepackage[]{microtype}
  \UseMicrotypeSet[protrusion]{basicmath} % disable protrusion for tt fonts
}{}
\makeatletter
\@ifundefined{KOMAClassName}{% if non-KOMA class
  \IfFileExists{parskip.sty}{%
    \usepackage{parskip}
  }{% else
    \setlength{\parindent}{0pt}
    \setlength{\parskip}{6pt plus 2pt minus 1pt}}
}{% if KOMA class
  \KOMAoptions{parskip=half}}
\makeatother
\usepackage{graphicx}
\makeatletter
\newsavebox\pandoc@box
\newcommand*\pandocbounded[1]{% scales image to fit in text height/width
  \sbox\pandoc@box{#1}%
  \Gscale@div\@tempa{\textheight}{\dimexpr\ht\pandoc@box+\dp\pandoc@box\relax}%
  \Gscale@div\@tempb{\linewidth}{\wd\pandoc@box}%
  \ifdim\@tempb\p@<\@tempa\p@\let\@tempa\@tempb\fi% select the smaller of both
  \ifdim\@tempa\p@<\p@\scalebox{\@tempa}{\usebox\pandoc@box}%
  \else\usebox{\pandoc@box}%
  \fi%
}
% Set default figure placement to htbp
\def\fps@figure{htbp}
\makeatother
\setlength{\emergencystretch}{3em} % prevent overfull lines
\providecommand{\tightlist}{%
  \setlength{\itemsep}{0pt}\setlength{\parskip}{0pt}}
\usepackage{amsmath,amssymb,mathtools,needspace}
\xeCJKDeclareCharClass{CJK}{"2032,"0400->"04FF,"2160->"216B,"2460->"2473,"25A0,"25C7,"25CB,"25CF}
\IfFontExistsTF{Arial Unicode MS}{\xeCJKsetup{AutoFallBack=true}\setCJKfallbackfamilyfont{\CJKrmdefault}{Arial Unicode MS}}{}
\setcounter{secnumdepth}{0}
\setlength{\emergencystretch}{2em}
\allowdisplaybreaks
\usepackage{bookmark}
\IfFileExists{xurl.sty}{\usepackage{xurl}}{} % add URL line breaks if available
\urlstyle{same}
\hypersetup{
  pdftitle={引言},
  colorlinks=true,
  linkcolor={Maroon},
  filecolor={Maroon},
  citecolor={Blue},
  urlcolor={Blue},
  pdfcreator={LaTeX via pandoc}}

\title{引言}
\author{}
\date{}

\begin{document}
\maketitle

\subsection{2.1 引言}\label{ux5f15ux8a00}

许多实际问题都要用函数 \(y=f(x)\) 来表示某种内在规律的数量关系，其中相当一部分函数是通过实验或观测得到的。虽然 \(f(x)\) 在某个区间 \([a,b]\) 上是存在的，有的还是连续的，但却只能给出 \([a,b]\) 上一系列点 \(x_i\) 的函数值 \(y_i=f(x_i)\)（\(i=0,1,\cdots,n\)），这只是一张函数表。有的函数虽有解析表达式，但由于计算复杂，使用不方便，通常也造一个函数表，如大家熟悉的三角函数表、对数表、平方根和立方根表等。为了研究函数的变化规律，往往需要求出不在表上的函数值。因此，我们希望根据给定的函数表构造一个既能反映函数 \(f(x)\) 的特性、又便于计算的简单函数 \(P(x)\)，用 \(P(x)\) 近似 \(f(x)\)。通常选一类较简单的函数如代数多项式或分段代数多项式作为 \(P(x)\)，并使 \(P(x_i)=f(x_i)\) 对于 \(i=0,1,\cdots,n\) 成立，这样确定的 \(P(x)\) 就是我们希望得到的插值函数。例如，在现代机械工业中用计算机程序控制加工机械零件，根据设计可给出零件外形曲线的某些型值点 \((x_i,y_i)\)（\(i=0,1,\cdots,n\)），加工时为控制每步走刀方向及步数，就要算出零件外形曲线其他点的函数值，才能加工出外表光滑的零件，这就是求插值函数的问题。下面给出有关插值法的定义。

\textbf{定义 2.1} 设函数 \(y=f(x)\) 在区间 \([a,b]\) 上有定义，且已知在点 \(a\le x_0\le x_1<\cdots<x_n\le b\) 上的值 \(y_0,y_1,\cdots,y_n\)，若存在一简单函数 \(P(x)\)，使

\[
P(x_i)=y_i\qquad (i=0,1,\cdots,n)
\tag{2.1.1}
\]

成立，就称 \(P(x)\) 为 \(f(x)\) 的\textbf{插值函数}。点 \(x_0,x_1,\cdots,x_n\) 称为\textbf{插值节点}，包含插值节点的区间 \([a,b]\) 称为\textbf{插值区间}，求插值函数 \(P(x)\) 的方法称为\textbf{插值法}。若 \(P(x)\) 是次数不超过 \(n\) 的代数多项式，即

\[
P(x)=a_0+a_1x+\cdots+a_nx^n,
\tag{2.1.2}
\]

其中 \(a_i\) 为实数，就称 \(P(x)\) 为\textbf{插值多项式}，相应的插值法称为\textbf{多项式插值}；若 \(P(x)\) 为分段的多项式，就称之为\textbf{分段插值}；若 \(P(x)\) 为三角多项式，就称之为\textbf{三角插值}。本章只讨论多项式插值与分段插值。

从几何图形上看，插值法就是求曲线 \(y=P(x)\)，使其通过给定的 \(n+1\) 个点 \((x_i,y_i)\)，\(i=0,1,\cdots,n\)，并用它近似已知曲线 \(y=f(x)\)，如图 2.1 所示。

插值法是一种古老的数学方法，它来自生产实践。早在一千多年前，我国科学家在研究历法中就应用了线性插值与二次插值，但它的基本理论和结果却是在微积分

\begin{quote}
本页末句续 PDF 28。
\end{quote}

\subsubsection{校注（agent 补充）：NA5E-E001}\label{ux6821ux6ce8agent-ux8865ux5145na5e-e001}

扫描原页中，定义 2.1 的前两个节点确实印为 \(x_0\le x_1\)；这里忠实保留。它与 2.2.1 证明中``\(i\ne j\) 时 \(x_i\ne x_j\)''的前提不一致。普通 Lagrange 插值存在唯一性要求节点互异，使用时应采用

\[
a\le x_0<x_1<\cdots<x_n\le b.
\]

核验反例：当 \(n=1\)、\(x_0=x_1=0\)、\(y_0=y_1=c\) 时，所有 \(P(x)=c+tx\)（任意实数 \(t\)）都满足节点条件，故不唯一。这里的严格不等号是有明确依据的校注，不能把它当作原书实际印字。

\href{../../source-images/pdf-028.jpeg}{原页：书页 15／PDF 28}

\begin{quote}
页首接 PDF 27 的``微积分''。
\end{quote}

产生以后才逐步完善的，其应用也日益增多。特别是在电子计算机广泛使用以后，由于航空、造船、精密机械加工等实际问题的需要，插值法在实践上和理论上显得更为重要，由此得到进一步发展，尤其是近几十年发展起来的样条（spline）插值，获得了更为广泛的应用。

本章主要研究如何求出插值函数（包括分段插值函数及样条插值函数），并讨论插值函数 \(P(x)\) 的存在唯一性、收敛性及误差估计等。

\includegraphics[width=0.85\linewidth,height=\textheight,keepaspectratio,alt={图 2.1（原图裁切）}]{../assets/ch02a/figure-2.1.png}

图 2.1（原图）：曲线 \(y=P(x)\) 与 \(y=f(x)\) 通过共同的插值节点。全部文字标签：\(y\)、\(x\)、\(O\)、\(x_1\)、\(x_n\)、\(y_1\)、\(y_n\)、\(y=f(x)\)、\(y=P(x)\)。

\subsection{本节覆盖原页的校注记录}\label{ux672cux8282ux8986ux76d6ux539fux9875ux7684ux6821ux6ce8ux8bb0ux5f55}

下列记录按原页关联，含同页相邻小节的校注；计算时同时读取上文校注和完整原页。

\begin{itemize}
\tightlist
\item
  \texttt{NA5E-E001}：原书 PDF 页 27
\end{itemize}

\end{document}
